\documentclass[11pt]{article}

% ================= Packages =================
\usepackage[a4paper,margin=2.5cm]{geometry}
\usepackage{setspace}
\usepackage{times}
\usepackage{amsmath,amssymb}
\usepackage{graphicx}
\usepackage{hyperref}
\usepackage{enumitem}
\usepackage{titlesec}


\usepackage{fancyhdr}

\pagestyle{fancy}
\fancyhf{}

\fancyfoot[R]{\footnotesize\url{https://ijdant-math.org/}}
\fancyfoot[C]{\thepage}

\renewcommand{\headrulewidth}{0pt}
\renewcommand{\footrulewidth}{0.4pt}





% ================= Formatting =================
\onehalfspacing
\setlength{\parindent}{0pt}
\setlength{\parskip}{6pt}

\titleformat{\section}{\bfseries\large}{\thesection.}{0.5em}{}
\titleformat{\subsection}{\bfseries}{\thesubsection.}{0.5em}{}

% ================= Custom Commands =================
\newcommand{\journal}{\textbf{Communications on Applied Nonlinear Analysis}}
\newcommand{\issn}{ISSN: 1074-133X}
\newcommand{\volume}{Vol. X No. Y (20--)}

% ================= Document =================
\begin{document}


% ================= Journal Header =================
\noindent
\begin{minipage}[t]{0.85\textwidth}
\textbf{International Journal of Diophantine Approximation and Number Theory}\\
ISSN: 3117-812X\\
Vol. X No. Y (20--)
\end{minipage}
\vspace{0.7cm}
\hrule
\vspace{0.7cm}


% ================= Title =================
{\Large{$Original\ Research$}}\\
%\vspace{0.2cm}
\\
% {\LARGE\bfseries Paper Title}\\[0.5cm]
{\LARGE\bfseries Obtaining Exponents in Fermat, Markov and Catalan Diophantine Equations}\\[0.5cm]






%\vspace{0.8cm}

% ===== Line before Abstract =====
%\hrule\vspace{0.4cm}

% ================= Abstract Block =================

\noindent
\textbf{Pawel Jan Piskorz}$^{1,*}$
% \textbf{Firstname Lastname}$^{2}$,
% \textbf{Firstname Lastname}$^{2,*}$

\vspace{0.2cm}

\noindent
\begin{minipage}[t]{0.32\textwidth}

\end{minipage}
\hfill
\begin{minipage}[t]{0.65\textwidth} 

\noindent
$^{1}$ Independent Researcher; \texttt{paweljanpiskorz@gmail.com}\\
% $^{2}$ Affiliation 2; \texttt{e-mail@e-mail.com}
% \\
\noindent
$^{*}$\textit{Correspondence:} \texttt{ul.~Krakowska 55, 31-066 Kraków, Poland}

\vspace{1cm}



\end{minipage}


\noindent
\begin{minipage}[t]{0.32\textwidth}
%\vfill
%\textbf{Article History}\\[4pt]
Received: dd-mm-yyyy\\
Revised: dd-mm-yyyy\\
Accepted: dd-mm-yyyy
\\
{\small Copyright:©2026 by the authors.
Submitted to the IJDANT for possible open access
publication.}
\end{minipage}
\hfill
\begin{minipage}[t]{0.65\textwidth} 

\textbf{Abstract}\\
We propose a procedure which allows to determine the only admissible exponents
in Diophantine equations of three kinds: Fermat, Markov and Catalan.
First, we use a set of rewriting rules for each Diophantine equation.
Second, we obtain the parametrized versions of the Diophantine equations. 
Next, we differentiate each of parametrized equations obtaining the identities
for Fermat, Markov and Catalan equations. 


\textbf{Objectives:} The objective is to find the only admissible exponents in Fermat, Markov
and Catalan Diophantine equations.

\textbf{Methods:} We obtain the parametrized versions of the Diophantine equations, introduce a real parameter
and differentiate the appropriate equations with respect to it.

\textbf{Results:} Due to the determined identities we are able
to compute the equations admissible exponents.

\textbf{Conclusions:} We present the appropriate and only valid exponents in the three Diophantine equations: Fermat, Markov and Catalan.

\vspace{1cm}
% ===== Line after Abstract =====
\hrule


\end{minipage}

\vspace{1cm}

\textbf{Keywords:} Diophantine equations of Fermat, Markov and Catalan, Partial differentiation.

\textbf{Mathematics subject classification:} 11D61.


\section{Introduction} \label{sec:Intr}

We start with a set of rewriting rules for the Fermat, Markov and Catalan equations
transforming those into equations which allow their partial differentiation.
This leads us to certain identities. From the obtained identities we
receive a sequence the limits of which determine the equations admissible exponents.


\section{Fermat Diophantine equation}

We can write the Fermat's Diophantine equation \cite{Edwards1978FLT, Vandiver1929FLT, Vandiver1946FLT, Dickson1917FLT} as
\begin{equation}
\label{eq:FDE}
a^{n} + b^{n} = c^{n},
\end{equation}
where $a, b, c, n$ are natural numbers (i.e.~integers greater than zero).


We will replace $a$ with $x + kx$, where $x \in \mathbb{Q}$ and $k \in \mathbb{N}$, rewriting (\ref{eq:FDE}) as
\begin{equation}
\label{eq:FDE_rewritten}
(x + kx)^{n} + b^{n} = c^{n}
\end{equation}
or
\begin{equation}
\label{eq:FDE_rewritten_x}
(k + 1)^{n} x^{n} + b^{n} = c^{n}.
\end{equation}

Setting $x = \frac{1}{k+1}, \frac{2}{k+1}, \frac{3}{k+1}, \ldots$,
will make $a = 1, 2, 3, \ldots$ , in (\ref{eq:FDE}).
Similarly, we will have
\begin{equation}
\label{eq:FDE_rewritten_Y}
a^{n} + (y + ky)^{n} = c^{n}
\end{equation}
and
\begin{equation}
\label{eq:FDE_rewritten_Z}
a^{n} + b^{n} = (z + kz)^{n},
\end{equation}
where $y, z = \frac{1}{k+1}, \frac{2}{k+1}, \frac{3}{k+1}, \ldots$ \@.

%%%%%%%%%%%%%%%%%%%%%%%%%%%%%%%%%%%%%%%%%%%%%%%%%%%%%%%%%%%%%%%%%

We introduce a real parameter
$p \in \mathbb{R}$ into (\ref{eq:FDE_rewritten}) as follows
\begin{equation}
\label{eq:FDE_rewritten_with_p}
(x + kpx)^{n} + b^{n} p = c^{n} p.
\end{equation}
Additional parameters $q, r \in \mathbb{R}$ will also be appropriately introduced into 
(\ref{eq:FDE_rewritten_Y})
and
(\ref{eq:FDE_rewritten_Z}) as
\begin{equation}
\label{eq:FDE_rewritten_Y_with_q}
a^{n} q + (y + kqy)^{n} = c^{n} q
\end{equation}
and
\begin{equation}
\label{eq:FDE_rewritten_Y_with_r}
a^{n} r + b^{n} r = (z + krz)^{n},
\end{equation}
respectively. Taking partial derivative of (\ref{eq:FDE_rewritten_with_p})
with respect to $p$, we obtain
\begin{equation}
\label{eq:FDE_rewritten_with_p_deriv}
nk(x + kpx)^{n-1} x + b^{n} = c^{n}
\end{equation}
or
\begin{equation}
\label{eq:FDE_rewritten_with_p_deriv_or}
nk(kp + 1)^{n-1} x^{n} + b^{n} = c^{n}.
\end{equation}
Setting $p=1$, we get
\begin{equation}
\label{eq:FDE_identity}
nk(k + 1)^{n-1} x^{n} + b^{n} = c^{n}.
\end{equation}
The coefficients of $x^{n}$ in (\ref{eq:FDE_identity}) and (\ref{eq:FDE_rewritten_x}) must be equal, hence
\begin{equation}
\label{eq:FDE_identity_a}
nk(k + 1)^{n-1} = (k + 1)^{n}.
\end{equation}
Therefrom, we obtain
\begin{equation}
\label{eq:FDE_identity_b}
nk(k + 1)^{n}/(k + 1) = (k + 1)^{n}
\end{equation}
and in this way we may write the equation for the values
of the exponent $n$ as a function of $k$
\begin{equation}
n = \frac{k+1}{k}.
\end{equation}

We can now determine the admissible values for $n$ as a function of $k = 1, 2, 3,$ $\ldots, \infty$, 
which will be $2/1, 3/2, 4/3, \ldots, 1$, respectively. 
Equations (\ref{eq:FDE_rewritten_Y_with_q}) and (\ref{eq:FDE_rewritten_Y_with_r}) are treated in a similar fashion
with the parameters $q$ and $r$ and partial differentiation.
We reject all $n$ that are not natural numbers, leaving us with $n = 1$ and $n = 2$ as the only acceptable exponent values in~(\ref{eq:FDE})\@.


%    Markov Diophantine equation

\section{Markov Diophantine equation}\label{MarkovEquation}


We write the Markov Diophantine equation \cite{Markov1879, Frobenius1913, Gyoda2022} as
\begin{equation}
\label{eq:MarkovDiophantineEqn}
a^{n} + b^{n} + c^{n} = 3\mathit{abc},
\end{equation}
where $a, b, c, n \in \mathbb{N}$.
Without loss of generality we rewrite (\ref{eq:MarkovDiophantineEqn}) as
\begin{equation}
(x + kx)^{n} + b^{n} + c^{n} = 3\mathit{abc}
\end{equation}
or
\begin{equation}
\label{eq:MDE_rewritten_x}
(k + 1)^{n} x^{n} + b^{n} + c^{n} = 3\mathit{abc},
\end{equation}
where $a = x + kx$, $k \in \mathbb{N}$.
We also have $x = \frac{1}{k+1}, \frac{2}{k+1}, \frac{3}{k+1}, \ldots$\@.
In this way $a = 1, 2, 3, \ldots$~, as assumed in the Markov Diophantine
equation~(\ref{eq:MarkovDiophantineEqn})\@.

After introducing the real parameter $p$,
we differentiate both sides of (\ref{eq:MDE_rewritten_with_p}) with respect to $p$
\begin{equation}
\label{eq:MDE_rewritten_with_p}
(x + kpx)^{n} + b^{n}p + c^{n}p = 3\mathit{abc}p
\end{equation}
and we establish the admissible exponents, as in the case of Fermat Diophantine equation.
We come to the
conclusion that the only possible exponents in the Markov Diophantine equation~(\ref{eq:MarkovDiophantineEqn})
are 1 and 2.

The one solution for $n = 1$ of $a + b + c = 3\mathit{abc}$ and $a, b, c > 0$~is
\begin{equation}
1 + 1 + 1 = 3 \cdot 1 \cdot 1 \cdot 1
\end{equation}

\begin{table}[!ht]
\caption{Example solutions of the Markov Diophantine equation for $n = 2$\@. Any triple 
$(a, b, c)$ satisfying the Markov Diophantine equation is called a Markov triple, and the numbers are Markov numbers.}
\label{tab:solutions}
\vspace{0.5cm}
\centering
\begin{tabular}{c c c}
\hline \hline \\ [-2.0ex]  % Space after first hline only
$a$ & $b$ & $c$ \\[0.6ex]  % extra vertical space before next \hline
\hline    \\[-1.5ex]  % extra vertical space before next \hline
1 & 1 & 1 \\
1 & 1 & 2 \\
1 & 2 & 5 \\
2 & 5 & 29 \\[0.6ex]  % extra vertical space before next \hline
\hline \hline
\end{tabular}
\end{table}
Some example solutions for $n = 2$ of $a^{2} + b^{2} + c^{2} = 3\mathit{abc}$ are presented in Table~\ref{tab:solutions}\@.
They are:
\begin{align*} 
1^{2} + 1^{2} + 1^{2} =  1 + 1 + 1 =  3 \cdot 1 \cdot 1 \cdot 1 &= 3         \\ \newline
1^{2} + 1^{2} + 2^{2} =  1 + 1 + 4 =  3 \cdot 1 \cdot 1 \cdot 2 &= 6         \\ \newline
1^{2} + 2^{2} + 5^{2} =  1 + 4 + 25 = 3 \cdot 1 \cdot 2 \cdot 5 &= 30        \\ \newline
2^{2} + 5^{2} + 29^{2} = 4 + 25 + 841 = 3 \cdot 2 \cdot 5 \cdot 29 &= 870.   
\end{align*} 


%%%%%%%%%%%%%%%%%%%%%%%%%%%%%%%%%%%%%%%%%%%%%%%



%    Catalan Diophantine equation

\section{Catalan Diophantine equation}\label{CatalanEquation}


We write the Catalan Diophantine 
equation~\cite{Catalan1844, mihalescu2004} as



\begin{equation}
\label{eq:CatalanDiophantineEqn}
a^{i} - b^{j} = 1,
\end{equation}
where $a, b, i, j \in \mathbb{Z}$ are integers greater than one.

Without loss of generality we rewrite (\ref{eq:CatalanDiophantineEqn}) as
\begin{equation}
(x + kx)^{i} - b^{j} = 1
\end{equation}
or
\begin{equation}
\label{eq:CatalanDE_rewritten_x}
(k+1)^{i}x^{i} - b^{j} = 1,
\end{equation}
where $a = x + kx$ and $k \in \mathbb N$.
We also have $x = \frac{2}{k+1}, \frac{3}{k+1}, \ldots$~,
thus allowing $a = 2, 3, \ldots$~, as assumed in the Catalan Diophantine
equation~(\ref{eq:CatalanDiophantineEqn})\@.



After introducing the real parameter $p$,
we differentiate both sides of (\ref{eq:CatalanDE_rewritten_with_p}) with respect to $p$
\begin{equation}
\label{eq:CatalanDE_rewritten_with_p}
(x + kpx)^{i} - b^{j}p = p
\end{equation}
and we determine the exponents.
We come to the conclusion that the only admissible exponent $i$ in the Catalan Diophantine equation~(\ref{eq:CatalanDiophantineEqn})
is equal to 2.



With $i = 2$, we obtain the following Catalan Diophantine equation
\begin{equation}
\label{eq:CatalanDE_with_X_power_2}
a^{2} - b^{j} = 1
\end{equation}
or
\begin{equation}
\label{eq:CatalanDE_with_X_power_2_modified}
b^{j} = a^{2} - 1.
\end{equation}
Since $b > 1$, we can set $b = 2$, obtaining
\begin{equation}
2^{j} = a^{2} - 1.
\end{equation}
One can verify that for $j=3$ we have $2^{3} = 8$ and $a = 3$\@.
Then we have the Catalan equation solution
\begin{equation}
a^{i} - b^{j} = 1
\end{equation}
$a = 3, i = 2, b = 2, j = 3$,
which gives
\begin{equation}
3^{2} - 2^{3} = 1.
\end{equation}

%%%%%%%%%%%%%%%%%%%%%%%%%%%%%%%%%%%%%%%%%%%%%%%

\section{Conclusion}


%++++++++++++++++++++++
Using a relatively straightforward set of rewriting ru\-les for the Diophantine equations of Fermat, Markov, and
Catalan, we were able to transform each into a form that allowed the introduction of rational parameters 
$x, y, z$ and real parameters $p, q, r$\@. When differentiated with respect to the latter, the resulting equations gave us a way to determine the admissible values for the exponent $n$ in the Fermat and Markov, and exponents $i, j$ in Catalan Diophantine equations.


%++++++++++++++++++++++



\bibliographystyle{vancouver}
\bibliography{FerMarCat_IJDANT_ver_007.bib}


\end{document}

